\section{System Design}
STEI-ZStat is a single-channel, three-electrode potentiostat interface based on the Analog Devices AD5941 analog front-end (AFE) and the Seeed Studio XIAO RP2040 microcontroller unit (MCU) (Fig.~\ref{fig:system_design}). The hardware architecture is derived from the design in \cite{vamos2025hunstat2} and optimized to fit within a footprint of 63.1 $\times$ 45.0~mm. The AD5941 integrates a potentiostatic control loop with a transimpedance input stage and is the basis for several reported potentiostats \cite{bautista2025helpstat, bill2023electrochemical, vamos2025hunstat2}. It supports impedance measurements from 0.015~Hz to 200~kHz and excitation voltages of up to 1.2~V$_{pp}$ \cite{devices2024ad5940}, although this study experimentally validates the instrument only from 0.2~Hz to 200~kHz. Compared with the base HunStat2 design, STEI-ZStat introduces three hardware modifications to improve EIS performance.

The first modification is a jumper-selectable $R_{CAL}$ resistor array with values of 200~$\Omega$, 4.7~k$\Omega$, and 10~k$\Omega$, along with an unpopulated footprint for alternative resistance values. During EIS measurements, the known $R_{CAL}$ impedance is measured separately to determine the unknown impedance \cite{ad5940_eis_github}. The selectable $R_{CAL}$ values allow the calibration resistor to match the expected impedance of the device under test more closely, as recommended by Analog Devices.

The second modification is an on-board dummy cell comprising a Randles circuit and a 10~k$\Omega$ (1\%) resistor for self-testing, selectable via jumpers. Although the on-board Randles circuit uses the same nominal component values as the external dummy cell shown in Fig.~\ref{fig:system_design}, the external cell was used for validation because it is populated with tighter-tolerance components ($R_s$: 560~$\Omega \pm$ 0.1\%, $R_{ct}$: 10~k$\Omega \pm$ 0.02\%, $C_{dl}$: 33~nF $\pm$ 5\%), providing a more reliable reference.

The third modification is the PCB layout, which uses the partitioned pinout of the AD5941 to separate analog and digital signal paths. The microcontroller and digital traces are physically isolated from sensitive analog nodes, including electrode connections, calibration resistors, and reference capacitors. This layout maintains a continuous ground plane and is intended to minimize electromagnetic interference (EMI) coupling into high-frequency EIS measurements.

The firmware is implemented in C/C++ using the Arduino IDE and is based on the Analog Devices AD5940/AD5941 library and example code \cite{ad5940_eis_github}, while the Python GUI uses Tkinter and Matplotlib. The software also supports exporting measurement data to a comma-separated values (CSV) file. In addition to EIS, cyclic voltammetry (CV), open-circuit potential (OCP), chronoamperometry (CA), square-wave voltammetry (SWV), and differential pulse voltammetry (DPV) have been implemented but have not yet been experimentally validated.

The unit build cost is US\$25.36 at a build quantity of five (May 2026). This includes PCB fabrication and assembly of the AFE and passive components (US\$20.34), the MCU (US\$4.80), screen-printed electrode (SPE) connectors (US\$0.09), and switches, jumpers, and pin headers (US\$0.13). The total excludes shipping, an enclosure, cables, electrodes, a host computer, and assembly labor. Cost and dimensional comparisons against previously reported systems are summarized in Table~\ref{tab:comparison}.